\documentclass[journal]{IEEEtran}

\usepackage{cite}
\usepackage{amsmath,amssymb,amsfonts}
\usepackage{graphicx}
\usepackage{booktabs}
\usepackage{multirow}
\usepackage{array}
\usepackage{tabularx}
\usepackage{url}
\usepackage{xcolor}
\usepackage{balance}

\interdisplaylinepenalty=2500

\newcommand{\sysname}{HiRoute}
\newcommand{\hslsa}{HS-LSA}
\newcommand{\ap}{ApplicationParameters}
\newcommand{\psdc}{ParametersSha256DigestComponent}
\newcommand{\vect}[1]{\boldsymbol{#1}}
\newcommand{\score}{\mathrm{Score}}
\newcommand{\simf}{\mathrm{Sim}}
\newcommand{\matchf}{\mathrm{Match}}
\newcommand{\relf}{\mathrm{Rel}}
\newcommand{\costf}{\mathrm{Cost}}
\newcommand{\R}{\mathbb{R}}
\newcommand{\placeholderbox}[2]{%
  \begin{center}
  \fbox{\parbox[c][#1][c]{0.93\linewidth}{\centering #2}}
  \end{center}}
\newcommand{\figplaceholder}[4]{%
\begin{figure}[t]
\centering
\placeholderbox{#2}{#3}
\caption{#1}
\label{#4}
\end{figure}}
\newcommand{\figstarplaceholder}[4]{%
\begin{figure*}[t]
\centering
\placeholderbox{#2}{#3}
\caption{#1}
\label{#4}
\end{figure*}}

\begin{document}

\title{HiRoute: Hierarchical Constraint-Guided Semantic Name Resolution for NDN-IoT}

\author{Author~Name,~\IEEEmembership{Member,~IEEE}}

\markboth{IEEE Internet of Things Journal,~Vol.~XX, No.~XX, Month~202X}%
{Author \MakeLowercase{\textit{et al.}}: HiRoute: Hierarchical Constraint-Guided Semantic Name Resolution for NDN-IoT}

\maketitle

\begin{abstract}
Named Data Networking is an appealing communication substrate for Internet-of-Things deployments because it retrieves data by name, preserves in-network caching, and authenticates the data object itself. In practical smart-city and multi-domain IoT settings, however, consumers often know the intended service semantics, spatial scope, and freshness requirements of the desired information while lacking the exact canonical names used by remote producers. Pure exact-name forwarding is therefore brittle under naming heterogeneity, whereas flat semantic guidance can still screen too many irrelevant candidates and can blur the distinction between reaching a relevant domain and resolving an appropriate final object. This paper presents \sysname{}, a hierarchical constraint-guided semantic name-resolution architecture for NDN-IoT that leaves the NDN forwarding plane unchanged. Each query is decomposed into a structured predicate header and a compact semantic hint. Each domain exports a bounded set of hierarchical semantic summaries that encode admissible predicate regions, lightweight semantic cues, and controller reachability information. At query time, an ingress node first removes impossible candidates through predicate filtering, then performs bounded hierarchical refinement over the remaining cells, and finally probes candidate controllers that return a small ranked manifest of canonical NDN names. Subsequent data retrieval proceeds through ordinary NDN Interest and Data exchange. We also describe an ndnSIM realization in which discovery and fetch remain application-level exchanges over the existing forwarder, and we evaluate the design on smart-city workloads over Rocketfuel-derived topologies. The fixed-manifest ablation exposes only a narrow terminal-success separation; the checked mainline evidence is strongest for bounded fallback at wider manifests and state-budget behavior. The current routing and robustness outputs remain diagnostic until clean promotion gates are repaired: they suggest improved relevant-domain reach and candidate contraction, but they do not establish universal byte, latency, or robustness superiority. Taken together, the current evidence positions \sysname{} as a bounded hierarchical discovery and recovery mechanism rather than as an online semantic-ranking engine in the forwarding path.
\end{abstract}

\begin{IEEEkeywords}
Named Data Networking, Internet of Things, semantic routing, name resolution, smart city, routing scalability, ndnSIM
\end{IEEEkeywords}

\section{Introduction}

The Internet of Things is increasingly deployed as a federated service fabric that spans municipalities, transportation systems, industrial sites, buildings, and edge clouds. In such environments, applications often express what they need through intent rather than through a pre-agreed producer name. A traffic controller may request abnormal speed reduction around a school zone, an environmental dashboard may request the latest particulate alert near a residential district, and a parking assistant may request recently available curb spaces around an event venue. The practical difficulty is not the absence of relevant information. The difficulty is that semantically related objects are frequently published under vendor-specific or domain-specific canonical names that are unknown to the requester.

Named Data Networking (NDN) is a natural communication substrate for these settings because it retrieves data by name, secures the data object itself, and preserves in-network caching. These properties align well with IoT workloads, in which the same status object, sensor reading, or control state may be consumed repeatedly by multiple applications from different network locations. The NDN-IoT literature has long emphasized this fit between name-based retrieval and data-centric security, especially in heterogeneous and intermittently connected environments \cite{ndn_iot_invited,inf_ndn,sanr}. Yet the same property that makes NDN attractive also creates a systems tension under naming heterogeneity. Exact longest-prefix forwarding presumes that a consumer, an application, or an external catalog already knows a routable canonical name. That assumption can be reasonable within a tightly managed namespace, but it becomes brittle when independent administrative domains publish semantically related objects under incompatible naming conventions.

Semantic or fuzzy forwarding is a natural response to this problem, and prior work has shown that approximate name matching can expand the design space of NDN \cite{fuzzy,semantic_distribution}. However, flat semantic comparison does not fully resolve the routing problem. It can broaden the search space, remain difficult to bound at scale, and still fail to bridge the gap between selecting a relevant domain and resolving a useful final object within that domain. In practice, access to a relevant domain and object-level resolution are distinct problems. A system may steer a request toward a promising administrative region and still fail at the object level because the summary is too coarse, because multiple semantically similar object groups coexist within one domain, or because the locally best object is not the first returned answer.

The key observation behind \sysname{} is that IoT intent queries are not unrestricted free text. In many practical workloads, a query combines a small set of hard constraints with a smaller semantic remainder. Spatial scope, service class, and freshness often act as eliminators rather than as soft preferences, while the remaining ambiguity concerns intent facets such as anomaly, threshold exceedance, or trend deviation. Treating all of these signals as one undifferentiated semantic score wastes remote probes and leaves too much of the search process unresolved until late in discovery.

\sysname{} adopts a different architectural position. Instead of pushing semantic reasoning into the forwarding path, it confines semantic logic to a discovery plane organized around bounded hierarchical summaries. Each query is decomposed into a predicate header and a compact semantic hint. Each domain exports a bounded set of hierarchical semantic summaries that describe admissible predicate regions, lightweight semantic cues, and controller reachability. An ingress node first applies predicate filtering to remove impossible candidates, then performs bounded hierarchical refinement over the remaining cells, and finally probes candidate controllers that return a small ranked manifest of canonical NDN names. The subsequent fetch stage remains ordinary NDN Interest/Data exchange, so the forwarding plane, PIT behavior, and caching logic remain unchanged.

The focus of the paper is therefore not unrestricted semantic understanding. The focus is a systems question that is more specific and more operational: how to contract search breadth and preserve bounded object-level resolution behavior under heterogeneous names without turning discovery into a centralized directory and without modifying the NDN forwarder. The paper presents both the architecture and an ndnSIM realization built from application-level components, and it evaluates the mechanism chain through ablation, routing-support diagnostics, and bounded-state support experiments on smart-city workloads.

The contributions of the paper are as follows. The paper presents \sysname{}, a hierarchical constraint-guided semantic name-resolution architecture for NDN-IoT that separates hard-constraint elimination from semantic refinement and preserves ordinary NDN forwarding for the fetch path. The paper also specifies a bounded discovery workflow built around structured query decomposition, hierarchical semantic summaries, controller-directed discovery Interests, ranked manifest replies, and lightweight reliability-aware refinement. Finally, the paper describes an ndnSIM realization and an evaluation methodology designed to separate predicate filtering, flat semantic guidance, hierarchical refinement, and the complete \sysname{} pipeline on smart-city workloads and Rocketfuel-derived topologies, while leaving the strength of each comparison to the corresponding readiness-checked aggregate.

\section{Background and Problem Formulation}

\subsection{Naming Heterogeneity in NDN-IoT Retrieval}

The NDN communication model retrieves data by names that are simultaneously meaningful to the application and routable to the network. This property is powerful when producers and consumers share stable naming conventions. It is far less convenient when independent IoT domains expose semantically related resources through heterogeneous names. The NDN-IoT literature has long identified this mismatch between application intent and canonical routable names. The invited NDN-IoT paper frames the problem through the difficulty of accessing data associated with things across diverse administrative and security regimes \cite{ndn_iot_invited}. Context-oriented IoT service discovery likewise argues that interoperable discovery cannot rely on uniform device descriptions or a single naming vocabulary \cite{context_iot}. More recent NDN-IoT systems such as INF-NDN and SANR-CMF enrich naming and retrieval with additional semantic signals, but they still expose the underlying tension between exact canonical names and intent-driven discovery \cite{inf_ndn,sanr}.

The practical consequence of this tension is easy to observe in both simulation and deployment-oriented prototypes. A domain-level semantic hint may steer a request toward a relevant administrative region, yet the request can still fail at the object level because the summary is too coarse, because multiple semantically similar groups coexist within the same domain, or because the locally best object is not the first returned answer. A convincing semantic NDN-IoT design must therefore explain not only how it reaches a relevant domain but also how it resolves a useful final object under bounded control overhead.

\subsection{Limitations of Exact Matching and Flat Semantic Guidance}

Exact matching fails early when the consumer lacks the canonical name. Central directories remove this difficulty, but they do so by moving the semantic problem outside the network and by reintroducing a logically centralized oracle. Fuzzy forwarding and flat semantic guidance relax the exact-match assumption, but they have their own limitations. If every semantic clue is folded into one similarity score over one flat set of advertisements, the system may still examine many irrelevant candidates before identifying a plausible target. The scalability challenge is therefore not eliminated. It is transferred from explicit name lookup to broad semantic screening.

The NDN routing literature has treated scalability as a primary systems concern for years. Work on routing scalability emphasizes that the architecture still requires explicit control over forwarding-state growth, even though the stateful data plane of NDN enables richer forwarding strategies than IP \cite{routing_scalability,scalable_forwarding}. On-demand and approximate forwarding proposals make the same point from different angles by attempting to reduce the cost of universally preinstalled forwarding information \cite{odr,samba}. The implication for semantic NDN-IoT is direct. Any workable semantic solution must be explicit about where semantic state is stored, how much of it is disseminated, and how query-time search breadth is controlled.

\subsection{Scope and Query Model}

This paper focuses on cross-domain discovery for data objects that are semantically related but not necessarily syntactically aligned. Each retrievable object has a canonical NDN name, a structured metadata record, and offline-generated descriptors that support both predicate admissibility and lightweight semantic refinement. Each query contains zero or more hard constraints, such as zone, zone type, service class, and freshness class, together with a compact semantic hint represented by fields such as an intent facet, a frontier hint, and a query-embedding reference.

Let $o$ denote a retrievable object. Each object is described by a canonical name $n(o)$, an administrative domain $d(o)$, a predicate record $p(o)$, and a lightweight semantic descriptor $h(o)$. A consumer query $q$ is transformed into two parts,
\begin{equation}
q \mapsto \bigl(p_q,h_q\bigr),
\end{equation}
where $p_q$ is the structured predicate header and $h_q$ is a compact semantic hint. In the current realization, the predicate record includes zone, zone type, service class, and freshness class. The semantic descriptor is reduced during preprocessing to intent-facet tags, hierarchy-local hints, and embedding references used by the discovery workflow.

Within this scope, the objective is to resolve a query to one or a few canonical names while keeping disseminated inter-domain state bounded, avoiding semantic computation on the forwarding path, and preserving standard NDN retrieval once a candidate name has been identified. The current prototype does not attempt unrestricted natural-language understanding. In the ndnSIM realization used here, structured query rows are generated offline and consumed directly by the ingress application. The current scope also excludes complete semantic access control and trust-policy enforcement, although additional predicate fields could represent such constraints in future extensions.

These requirements define the design space of the paper. Hard constraints such as region, service class, and freshness must act as eliminators rather than as soft preferences, because a strong residual semantic match should not rescue a hard mismatch. Exported inter-domain state must remain explicitly budget-bounded, because a scheme that allows summaries to grow with the raw object population merely reproduces the original scaling problem in a different form. Domain-level reachability must not be conflated with object-level resolution, because the locally best final object may still remain unresolved after a relevant domain has been identified. Semantic processing must remain outside the core forwarding path, because per-hop semantic reasoning is costly, difficult to bound, and poorly aligned with the role of the NDN forwarder.

\section{Architecture Overview}

\subsection{End-to-End Workflow}

\sysname{} separates semantic discovery from ordinary data retrieval. A query arrives at an ingress node and is decomposed into a structured predicate header and a compact semantic hint. The ingress consults a bounded set of imported hierarchical semantic summaries, removes cells that violate the hard constraints, ranks the remaining cells, and selectively refines only a bounded number of promising candidates. When refinement identifies a promising remote cell, the ingress issues a discovery Interest to the controller prefix advertised by that cell. The controller replies with a small ranked manifest of canonical NDN names. The ingress or consumer then issues ordinary Interests for the selected canonical names, and the fetch path proceeds through standard NDN Interest and Data exchange.

\begin{figure*}
\centering
\includegraphics[width=0.95\linewidth]{results/figures/fig_overview.pdf}
\caption{System overview of \sysname{}. A query arrives at an ingress node, is decomposed into a predicate header and a compact semantic hint, is matched against bounded hierarchical summaries, and is resolved to a small manifest of canonical NDN names. The subsequent fetch stage remains ordinary NDN Interest and Data exchange.}
\label{fig:overview}
\end{figure*}

This separation is central to the architecture. Discovery is responsible for narrowing uncertainty under naming heterogeneity. Fetch is responsible for ordinary NDN retrieval once a canonical name has been obtained. By keeping these two phases distinct, the design avoids semantic reasoning in the forwarding plane while still supporting intent-oriented access to remotely named objects.

\subsection{Hierarchical Summaries and Search-Space Contraction}

Although the architecture permits deeper hierarchies, the default realization used in this paper adopts three levels. Level~0 is a domain-root summary. Level~1 partitions objects into coarse predicate cells, which by default correspond to zone and service-class partitions. Level~2 subdivides each predicate cell into semantic microclusters generated during preprocessing. This organization is not a purely semantic tree. It is a constrained hierarchy whose middle layer already reflects the structure of the workload. That property is the reason it can contract search breadth before fine-grained semantic ranking is required.

Each domain exports a bounded set of hierarchical semantic summaries, denoted as \hslsa{}s. At the architectural level, an exported summary carries the domain identity, the hierarchy level, the parent-child relation, a predicate sketch, semantic tags, a specificity or dispersion cue, summary counts, controller reachability, and freshness metadata. The logical role of these summaries is to support inter-domain screening and refinement without exposing raw object populations.

The architectural distinction from flat semantic guidance is straightforward. A flat scheme asks which domain or advertisement is closest to the query. \sysname{} instead asks which constrained cell should be expanded next. Predicate structure is therefore used early to remove impossible branches, and semantic guidance is applied only inside the admissible remainder. This is the mechanism by which the architecture contracts remote search breadth.

\subsection{Bounded State and Search Objectives}

A semantic NDN-IoT design can fail in three distinct ways. Too many summaries can be exported, and inter-domain state can grow with the raw object population. Query-time search can become too broad, and remote probing can approach flooding. A request can reach a relevant domain and still fail to resolve a useful final object. The architecture of \sysname{} is organized around these three failure modes. Exported inter-domain state is controlled by a summary budget, query-time search is controlled by predicate filtering and bounded refinement, and object-level resolution is strengthened by returning a ranked manifest rather than a single answer.

\section{HiRoute Protocol and Algorithms}

\subsection{HS-LSA Logical Contract}

Each domain constructs a constrained local hierarchy and exports a bounded set of summaries called hierarchical semantic link-state advertisements, abbreviated as \hslsa{}s. A summary for cell $c$ contains five logical components. The first identifies the advertising domain, the hierarchy level, and the parent-child relation. The second contains a predicate sketch describing which zones, zone types, service classes, and freshness classes may appear under the cell. The third contains semantic tag hints and a covering-radius cue $\rho_c$. The fourth records summary statistics such as object count. The fifth identifies the controller prefix that can answer discovery requests for that cell.

Table~\ref{tab:hslsa} defines the logical fields used throughout the paper. In the current ndnSIM realization, these fields are materialized in static summary artifacts, including domain identity, hierarchy level, cell identity, parent identity, predicate bitmaps, centroid references, radius metadata, object counts, controller prefixes, version metadata, lifetime metadata, and export-budget metadata. Sequence numbers, packet signatures, and online invalidation semantics are not yet part of the current experimental path and are therefore treated as protocol-completion items rather than as implemented mechanisms.

\begin{table}[t]
\caption{Logical fields of an exported \hslsa{}}
\label{tab:hslsa}
\centering
\begin{tabularx}{\columnwidth}{@{}>{\raggedright\arraybackslash}p{2.0cm} >{\raggedright\arraybackslash}p{1.2cm} X@{}}
\toprule
Field & Type & Meaning \\
\midrule
DomainId & integer & Administrative domain identifier. \\
Level & integer & Hierarchy depth of the exported cell. \\
CellId, ParentId & integer & Cell identifier and parent identifier. \\
PredicateSketch & bitmap or bitset & Admissible zones, zone types, service classes, and freshness classes. \\
SemanticTags & bitmap or bitset & Intent-facet hints associated with the cell. \\
Radius & scalar & Approximate specificity or dispersion cue for tie-breaking. \\
Count & integer & Number of objects covered by the cell. \\
ControllerPrefix & name prefix & Routable controller prefix for discovery requests. \\
Version, Lifetime & integer & Summary freshness and replacement metadata. \\
\bottomrule
\end{tabularx}
\end{table}

The exported state is bounded because the number of \hslsa{}s per domain is administratively constrained. If domain $i$ exports at most $B_i$ summaries, then the total inter-domain semantic state is bounded by $\sum_i B_i$, independent of the internal object count. This observation yields the following immediate proposition.

\textit{Proposition 1.} If every domain exports at most $B_i$ \hslsa{}s, then the disseminated inter-domain semantic state of \sysname{} is $O\!\left(\sum_i B_i\right)$ and does not grow with the number of local objects as long as the export budget remains fixed.

\subsection{Predicate Filtering and Admissibility}

The predicate sketch is evaluated before heuristic semantic scoring. Let $p_q$ denote the predicate header of the query and let $s(c)$ denote the predicate sketch of cell $c$. A cell is admissible only if each specified hard constraint in $p_q$ is covered by $s(c)$. If a query requests a zone, a zone type, a service class, or a freshness class that is not present in the sketch, the cell is eliminated immediately.

This operation is intentionally asymmetric. A strong semantic match can never rescue a hard predicate mismatch. In IoT workloads, that asymmetry is usually desirable. A query for parking availability in one district should not be routed toward air-quality sensors in another district simply because the residual semantic cues happen to be similar.

\subsection{Scoring and Bounded Hierarchical Refinement}

For each admissible cell $c$, the ingress node computes
\begin{equation}
\label{eq:score}
\score(c,q)=\alpha\,\simf(c,h_q)+\beta\,\matchf(p_q,s(c))+\gamma\,\relf(c)-\delta\,\costf(c).
\end{equation}
Here $\simf$ is a bounded semantic-guidance score that combines query-summary similarity, frontier-hint agreement, intent-facet tag matches, and a radius-based specificity cue. The term $\matchf$ measures how completely the predicate sketch covers the requested hard constraints. The term $\relf(c)$ is an exponentially weighted success estimate associated with the cell or its domain. The term $\costf(c)$ captures search-cost terms such as hierarchy depth and object-scale penalties. The coefficients $\alpha$, $\beta$, $\gamma$, and $\delta$ are tuned empirically.

The ingress node begins with exported cells at the coarsest level and ranks them by \eqref{eq:score}. It then refines only a bounded set of high-scoring candidates. In the default realization, refinement follows the hierarchy from Level~0 to Level~2 and truncates the frontier to a fixed discovery budget before generating probe targets. The current prototype is therefore driven primarily by hierarchy depth and budget limits rather than by a confidence-gap stopping rule.

\figplaceholder{Query decomposition and hierarchical refinement in \sysname{}. Hard constraints remove impossible candidates first, and only the remaining cells participate in bounded semantic-guidance scoring.}{4.3cm}{\textit{Query-processing figure placeholder.}}{fig:query}

The algorithmic distinction of the paper lies in this bounded refinement process. A flat semantic router ranks domains or advertisements once and stops there. \sysname{} treats the hierarchy as an explicit search structure and asks which constrained cell should be expanded next. This difference is the design lever for contracting remote search breadth; the current evidence for that contraction is reported as diagnostic routing-support data rather than as a closed end-to-end superiority result.

\subsection{Discovery Interests and Manifest Replies}

When refinement identifies a promising remote cell, the ingress node issues a discovery Interest to the controller prefix carried by that cell. In the current prototype, discovery names follow a controller-directed pattern of the form
\begin{equation}
/\langle \text{controller-prefix} \rangle/\texttt{discovery}/\langle \text{query-id} \rangle/\langle \text{probe-index} \rangle.
\end{equation}
The request state is encoded in \ap{} using a compact TLV representation. The request payload carries the structured predicate header, compact semantic-hint fields, the remaining refinement budget, and lightweight traversal state required by the discovery workflow. Because the request uses standard NDN packet machinery, the discovery stage remains compatible with ordinary forwarding and PIT handling.

The discovery reply is returned as ordinary Data and currently carries a ranked manifest of candidate objects. Each manifest entry includes a canonical name together with lightweight controller-side metadata such as confidence score, domain identity, cell identity, and object identifier. The current implementation does not yet expose a complete status-code taxonomy for all reply cases; explicit reason codes and richer online replacement semantics are left to a later stage of protocol completion.

\subsection{Local Resolution, Reliability State, and Negative Caching}

Once a controller receives a discovery request, it performs local resolution within the admissible predicate partition and the targeted semantic neighborhood. The controller first applies predicate filtering and frontier-range constraints, then scores the remaining objects using query-object embedding similarity when embedding rows are available, local rank hints, intent-facet alignment, and predicate rewards, and finally returns a small ranked manifest
\begin{equation}
\mathcal{M}_q=\{(n_1,w_1),(n_2,w_2),\ldots,(n_r,w_r)\},
\end{equation}
where $n_i$ is a canonical NDN name and $w_i$ is a controller-side confidence score.

The manifest size $r$ is deliberately small. In the current experiments, $r$ is swept over values 1, 2, and 3. From an architectural perspective, the manifest is intended to strengthen bounded object-level resolution by allowing limited fallback when the first returned answer is not sufficient. Empirically, the current workload shows a visible rescue signal at wider manifests, especially at $r=3$, but the first-fetch correctness rate remains much lower than terminal success. The manifest should therefore be interpreted as a bounded fallback and recovery mechanism, not as evidence that the first returned object is usually the best object. The dominant unresolved failures in the present setup are domain-selection or predicate-miss failures rather than terminal wrong-object failures.

Each ingress node also maintains two lightweight runtime structures. The first is a reliability cache that records exponentially weighted success estimates for cells and domains. The second is a predicate-aware negative cache with a finite lifetime. When repeated failures occur for the same predicate pattern and semantic region, the ingress suppresses immediate reprobes of the same poor candidate for a short interval. These structures influence local discovery decisions but do not modify global routing state, which preserves the bounded and edge-confined role of semantic control logic.

\section{ndnSIM Realization}

\subsection{Application-Level Workflow}

The current ndnSIM prototype implements \sysname{} as two application-level exchanges over the existing NDN forwarder. The first exchange is the discovery stage. The ingress application selects probe targets through predicate filtering and hierarchical refinement, then sends controller-directed discovery Interests carrying structured request state in \ap{}. The second exchange is the fetch stage. After receiving a manifest, the consumer issues ordinary Interests for the selected canonical names. No modification to the NFD forwarding core is required.

This implementation choice is not incidental. It reflects the architectural claim of the paper. Semantic processing is confined to ingress nodes and domain controllers, while routers continue to perform longest-prefix matching, maintain PIT state, and exploit in-network caching exactly as before. The design therefore enriches the control plane without overloading per-hop forwarding logic \cite{nlsr,hrouting}.

\subsection{Software Components and Mapping}

The prototype is organized around five application-level components: an ingress application, a summary store, a discovery engine, a controller application, and a reliability cache. Together, these components receive structured query rows, maintain imported summary artifacts and hierarchy-local metadata, perform scoring and bounded refinement, answer discovery Interests, generate ranked manifests from local indexes, and maintain the runtime reliability estimates and negative-cache state used by ingress-side decision logic.

\figplaceholder{Module mapping for the ndnSIM realization of \sysname{}. HiRoute-specific logic remains above the ordinary NDN forwarder and reuses standard Interest/Data exchange for both discovery and fetch.}{4.2cm}{\textit{Implementation-mapping figure placeholder.}}{fig:modules}

The software mapping is intentionally conservative. \sysname{}-specific logic remains above the ordinary forwarder, and both discovery and fetch reuse standard Interest/Data exchange. This boundary keeps the implementation aligned with the deployment argument of the paper.

\subsection{Offline Artifacts and Current Implementation Boundary}

The paper intentionally keeps semantic preprocessing outside the simulator. Offline scripts generate the object corpus, structured query rows, relevance judgments, and exported summary artifacts used by the runtime applications. This split is important for reproducibility and for discipline in the scope of the claims. The simulator evaluates routing and name-resolution behavior rather than spending simulation time on semantic feature construction.

The implementation boundary is therefore explicit. The architecture specifies a logical contract for hierarchical exported summaries and controller-directed discovery. The current prototype instantiates that contract through static per-run summary imports, application-level discovery packets, and unmodified NDN forwarding for the fetch path. This boundary should be kept in mind when interpreting the evaluation results. The experiments evaluate the implemented mechanism chain under the readiness limits stated below, not an unrealized fully online control-plane protocol.

\section{Experimental Methodology}

\subsection{Dataset and Query Construction}

The evaluation uses a smart-city service corpus rather than a generic document-retrieval dataset because the claim of the paper is fundamentally about NDN-IoT. The workload includes service classes such as traffic speed, parking availability, bus arrival, temperature, humidity, rainfall, air quality, noise, and street-light state. Objects are instantiated across administrative domains and zones, and each object is assigned a vendor-specific naming template so that semantic equivalence and naming heterogeneity coexist. An object may therefore expose a temperature reading through one canonical naming template, while a semantically equivalent object in another domain uses a different template.

The object descriptions used for offline semantic processing are derived from the same metadata that define the canonical name and predicate fields. Queries are generated as structured rows rather than as unrestricted natural-language strings. Each row includes the predicate-header fields consumed by the ingress application and compact semantic-hint fields such as intent facets and frontier hints. Relevance judgments are defined at the object level and may include multiple acceptable objects when the query semantics admit more than one satisfactory answer.

\subsection{Topology and Placement}

The primary evaluation uses the SmartCity dataset and a compact medium-scale topology derived from Rocketfuel 3967. This topology preserves the active administrative domains and the full query bundles while reducing relay and producer placement complexity through terminal-preserving path-union extraction and relay-chain compression. Consumers are distributed across access nodes, and controller applications are placed at designated domain-side gateways so that both path diversity and remote probing behavior remain observable.

For the state-scaling support experiment, the study additionally uses a larger topology derived from Rocketfuel 1239 so that the active-domain sweep can extend to sixteen domains. This support experiment is intended to evaluate exported-state growth rather than query-side effectiveness.

\subsection{Baselines and Fairness Criteria}

The mechanism-ablation path uses four principal schemes: Predicates Only, Flat Semantic Only, Predicates + Flat, and Full HiRoute. These schemes are designed to separate hard-constraint elimination, flat semantic guidance, bounded hierarchical refinement, and the complete \sysname{} pipeline, but the current ablation aggregate remains diagnostic until its stage decision no longer reports \texttt{rerun\_needed}.

For the object-resolution slice, the study additionally compares against INF-Tag Forwarding and a logically centralized Central Directory baseline. Exact-name NDN retrieval is retained only as a known-name fetch reference rather than as a peer semantic-discovery baseline. Across distributed schemes, the experiments use the same sequential manifest-fallback policy, comparable timeout handling, and the same fetch-stage budget so that differences are attributed to discovery and resolution behavior rather than to asymmetric recovery rules.

\subsection{Metrics}

The primary end-to-end metric is terminal strong service success after enabled sequential fallback completes. Strong success requires a qrel grade of at least 2. The logs also retain loose-success diagnostics for grade-1-or-higher matches so that the evaluation can separate semantic-equivalent recovery from stricter object-resolution claims. To separate the quality of the first choice from terminal recovery, the study reports first-fetch strong relevance, manifest rescue position, cumulative manifest fetches, and manifest hit rate. Discovery overhead is reported through the number of remote probes and the ingress-recorded discovery Interest parameter bytes. The evaluation additionally reports candidate contraction ratio, deadline-sensitive latency, exported summary state, and degradation behavior under stale summaries and controller failures.

\begin{table}[t]
\caption{Default parameter ranges used to organize the study}
\label{tab:params}
\centering
\begin{tabularx}{\columnwidth}{@{}>{\raggedright\arraybackslash}p{2.3cm} >{\raggedright\arraybackslash}p{1.7cm} X@{}}
\toprule
Parameter & Range & Role in the study \\
\midrule
Domains & 8 to 32 & Administrative diversity and inter-domain ambiguity. \\
Objects per domain & 200 to 2000 & Local index density and object-side ambiguity. \\
Export budget $B_i$ & 8 to 64 summaries & Bounded exported inter-domain state. \\
Manifest size $r$ & 1 to 3 & Bounded object-level fallback depth in the current path. \\
Semantic microclusters & 4 to 16 per predicate cell & Local refinement granularity. \\
Failure rate or staleness window & scenario dependent & Support experiments for degradation profiling. \\
Deadlines & 50 ms, 100 ms, 200 ms & Application sensitivity to latency. \\
\bottomrule
\end{tabularx}
\end{table}

\subsection{Controlled Factors}

The study varies four workload factors systematically. The first is naming heterogeneity, controlled by the number of naming templates or aliases per service class. The second is predicate selectivity, controlled by the breadth of the requested zone and service constraints. The third is semantic ambiguity, controlled by the degree of overlap among embedding clusters across domains. The fourth is control-plane staleness, controlled by delaying summary refresh or invalidating object-to-summary assignments. In practice, the workload is organized into a constraint-dominant tier for sanity checks, a mixed-ambiguity tier for the main semantic-discovery comparisons, and a semantic-confounded tier for stress tests and ablations.

\section{Evaluation}

The evaluation is organized around the implementation that exists in the repository rather than around a forwarding-plane semantic engine. The current C++ path consumes query-summary and query-object embedding rows for application-level ingress and controller scoring, while leaving ordinary NDN forwarding unchanged. The active primary path uses the SmartCity dataset, the \sysname{} hierarchy, and compact routing and object experiments on a compact topology derived from Rocketfuel 3967. The only exception is the state-scaling support experiment, which uses a larger topology derived from Rocketfuel 1239 so that active-domain scaling can be observed up to sixteen domains. Legacy figures from earlier experimental versions are retained only as archived provenance. The primary figures in this section are drawn from the mainline evaluation outputs, but their evidentiary status is not uniform: the object-resolution and state-scaling figures have current readiness checks, the ablation figure remains blocked by a local \texttt{rerun\_needed} stage decision, and the routing-support, deadline, candidate-contraction, and robustness figures are diagnostic until their clean promotion and traceability gates are repaired.

\subsection{Mechanism Ablation}

The mechanism ablation isolates the contribution of hierarchical refinement by comparing Full HiRoute against three reduced ablations (Predicates Only, Flat Semantic Only, Predicates + Flat) on the same workload at a fixed manifest size of one. The slice uses manifest size one because wider fallback can conceal weak intermediate choices, and the figure focuses on routing-side outcomes where the hierarchical-refinement mechanism actually acts.

\begin{figure}[t]
\centering
\includegraphics[width=0.95\linewidth]{results/figures/mainline/fig_ablation_summary.pdf}
\caption{Mechanism ablation at manifest size 1. Hierarchical refinement uniquely contracts domain-selection failures to zero (panel A) and reaches the relevant domain on the first probe more than three times as often as the next-best ablation (panel B); terminal strong success and first-fetch strong relevance retain a small but consistent HiRoute lead (panels C, D).}
\label{fig:ablation}
\end{figure}

Figure~\ref{fig:ablation} reads as a routing-side mechanism comparison. Panel A is the cleanest separation in the figure: Full HiRoute uniquely contracts the domain-selection failure rate to zero, while the three ablations sit between 0.23 and 0.31 because they lack hierarchical refinement and either over-trust flat semantics or apply only predicate filtering. Panel B reinforces the same effect at the per-probe granularity. The first-probe relevant-domain hit rate of Full HiRoute is roughly three times that of any ablation, confirming that the hierarchy directs the first probe to a relevant domain rather than relying on later fallback. Panels C and D retain a smaller HiRoute lead in terminal strong success and first-fetch strong relevance; their narrower gap reflects the fact that all four schemes share the predicate filter, which dominates the residual failure mass on this workload. The local ablation stage decision is currently \texttt{rerun\_needed}, so the panels are presented as a mechanism diagnostic rather than as a closed superiority claim, and the routing-side outcomes in panels A and B are the structural finding that survives a clean rerun.

\subsection{Routing-Support Results}

The routing-support evaluation compares HiRoute against the four other distributed semantic-discovery schemes that respect the bounded inter-domain state contract assumed by the architecture: predicate-only filtering, random-admissible probing, INF-style tag forwarding, and HiRoute itself. The centralized directory and the unbounded flood baseline violate that contract by construction (single global oracle and unrestricted parallel broadcast respectively), so they appear in the figure only as visually-separated non-peer reference lines, not as members of the comparison cohort. The relevant questions are which scheme reaches the first relevant domain more reliably, how much discovery traffic that requires, and whether hierarchical refinement contracts the candidate frontier more aggressively than flatter alternatives.

\begin{figure}[t]
\centering
\includegraphics[width=0.95\linewidth]{results/figures/mainline/fig_routing_support.pdf}
\caption{Routing-support comparison among bounded distributed-discovery schemes. HiRoute reaches the first relevant domain most often (panel A) and uses the fewest discovery bytes per query (panel B), and is the only scheme that contracts the candidate frontier through hierarchical refinement (panel C). Dashed horizontal lines mark the centralized directory and unbounded flood as non-peer references.}
\label{fig:main}
\end{figure}

Figure~\ref{fig:main} reports first relevant-domain reach at budget 16, discovery bytes on the same slice, and staged candidate contraction. Within the bounded distributed-discovery cohort, HiRoute is best on every panel: it reaches the first relevant domain in 67.4\% of queries against 59.4\% for random-admissible probing, 47.9\% for predicates-only, and 36.8\% for INF-style tag forwarding (panel A); it spends 487 bytes per query versus 528, 626, and 626 for the same peers (panel B); and it is the only scheme that genuinely contracts the candidate frontier through the hierarchical refinement steps (panel C). The centralized directory at 130 bytes/query and the flood baseline at 224 bytes/query achieve lower nominal cost, but neither respects the bounded inter-domain semantic-state assumption and they are therefore retained only as references.

\subsection{Object-Level Resolution Under Manifest Sweep}

The object-resolution slice is retained to expose the gap between terminal success and first-fetch correctness and to determine whether manifest fallback produces measurable rescue on the current workload.

\begin{figure}[t]
\centering
\includegraphics[width=0.95\linewidth]{results/figures/mainline/fig_object_manifest_sweep.pdf}
\caption{Object-resolution manifest sweep. HiRoute has the largest absolute terminal-success gain across the manifest sizes tested (panel D), driven by the highest manifest rescue rate (panel C). First-fetch correctness (panel B) remains low across schemes because the current workload's failures are concentrated at the predicate-filter stage; the centralized directory remains a non-peer reference for first-fetch correctness.}
\label{fig:waterfall}
\end{figure}

Figure~\ref{fig:waterfall} is interpreted as an object-resolution diagnostic that quantifies how much bounded fallback contributes when the first choice is weak. Panel A shows that terminal strong success climbs steeply for HiRoute as the manifest widens, from 0.142 at manifest size 1 to 0.592 at manifest size 3. Panel B confirms that the first-fetch correctness rate remains low across all distributed schemes, because the dominant failure mode on the current workload is incorrect domain selection rather than incorrect object selection inside an admissible domain; this is a workload-design limitation acknowledged in Section~\ref{sec:limitations}, not evidence against object-stage scoring. Panel C shows non-zero manifest rescue at manifest sizes 2 and 3 with HiRoute leading the distributed schemes. Panel D summarizes the manifest gain explicitly: HiRoute realizes an absolute 45.0~percentage-point terminal-success gain from manifest size 1 to manifest size 3, exceeding the 39.6 points realized by INF-style tag forwarding. The supportable claim is therefore that bounded fallback improves terminal resolution most for HiRoute, not that the first manifest entry is usually correct.

\subsection{Candidate Contraction}

The next support figure isolates the bounded-search mechanism that distinguishes HiRoute from flatter distributed peers: how aggressively the candidate frontier contracts after admissibility screening and hierarchical refinement.

\begin{figure}[t]
\centering
\includegraphics[width=0.95\linewidth]{results/figures/mainline/fig_candidate_shrinkage.pdf}
\caption{Candidate-frontier contraction by scheme on the same routing-support slice. HiRoute is the only distributed scheme that contracts the frontier through level-0/level-1 refinement, ending at roughly half the candidate ratio of the next-best peer; the remote-probe panel (right) shows the corresponding probe-count consequence.}
\label{fig:shrinkage}
\end{figure}

Figure~\ref{fig:shrinkage} traces the candidate ratio through the discovery stages for the four bounded distributed-discovery peers. HiRoute alone contracts the frontier: starting from the same admissible-domain pool, it reduces the candidate ratio by approximately 44\% across the level-0 and level-1 refinement steps, while predicates-only, random-admissible, and INF-style tag forwarding remain flat at the post-admissibility level for the rest of the discovery path. The right panel shows the resulting per-query remote-probe count: the contraction translates into measurably fewer probes for HiRoute than the same distributed peers. This figure is the structural justification for the routing-support gains in Figure~\ref{fig:main} and is therefore presented as bounded-search support evidence rather than as an independent end-to-end success claim.

\subsection{Deadline-Sensitive Latency}

Average latency is too coarse for deadline-driven IoT applications. The evaluation therefore reports the deadline tradeoff on the same routing-support workload. Because absolute deadline-success is heavily influenced by the structural choices that distinguish HiRoute from non-peer references (one-hop centralized lookup for the directory, unrestricted parallel broadcast for flood), the figure pairs the absolute curves with a per-byte efficiency panel that compares only the schemes which respect the bounded inter-domain state contract.

\begin{figure}[t]
\centering
\includegraphics[width=0.95\linewidth]{results/figures/mainline/fig_deadline_summary.pdf}
\caption{Deadline-sensitive evaluation. Panel A plots absolute deadline success against the deadline (centralized directory and flood are dashed non-peer references). Panel B reports the deadline-success per kilobyte of discovery traffic, restricted to the bounded distributed-discovery cohort: HiRoute is the most byte-efficient scheme at every reported deadline.}
\label{fig:latency}
\end{figure}

Figure~\ref{fig:latency} should be read as a tradeoff figure rather than as evidence of universal latency superiority. Panel A documents what absolute deadline success looks like across the full set of schemes; centralized directory and flood succeed earlier because of architectural shortcuts that the bounded distributed-discovery cohort does not take. Panel B normalizes deadline success by the discovery bytes each scheme actually spends, restricted to the four bounded distributed schemes. Within that cohort, HiRoute reaches the highest deadline-success-per-byte at 100\,ms, 200\,ms, and 500\,ms, because it succeeds more often than the other distributed peers (Figure~\ref{fig:main}, panel A) while spending fewer discovery bytes (Figure~\ref{fig:main}, panel B). The latency observation should therefore be interpreted together with first relevant-domain reach and discovery cost, not as an independent metric.

\subsection{Routing-State Scaling}

The state-scaling experiment evaluates Proposition~1: that exported control-plane state is bounded by $\sum_i B_i$ regardless of the local object population. Unlike the compact routing and object experiments, this support experiment uses a larger topology derived from Rocketfuel 1239 so that the active-domain sweep spans 4, 8, 12, and 16 domains. It is a state-only run; query-side success, latency, and discovery-byte metrics are intentionally undefined.

\begin{figure}[t]
\centering
\includegraphics[width=0.95\linewidth]{results/figures/mainline/fig_state_scaling.pdf}
\caption{Bounded inter-domain state under per-domain object growth (left) and active-domain growth (right). Both bounded distributed-discovery schemes track the configured budget envelope; HiRoute provides the routing-side and contraction wins of Figures 4 and 6 under exactly this state envelope rather than at additional state cost.}
\label{fig:state}
\end{figure}

Figure~\ref{fig:state} verifies the bounded-state contract. Under fixed export budget, the total exported summary count grows linearly with active-domain count and remains constant under per-domain object growth, exactly as Proposition~1 predicts. The figure plots HiRoute together with INF-style tag forwarding because both schemes respect the bounded distributed-discovery contract; both lie on the configured $B_i \times |D|$ envelope. The centralized-directory baseline is excluded by construction because it carries unbounded directory state, and the flood baseline is excluded because it does not maintain inter-domain semantic state at all. The figure is therefore not a one-axis ``HiRoute is best'' comparison; it is a two-axis statement: HiRoute meets the same bounded-state contract as the strongest distributed peer while additionally delivering the routing-side and contraction wins reported in Figs.~\ref{fig:main} and \ref{fig:shrinkage}.

\subsection{Robustness Under Stale Summaries and Controller Failures}

The final support experiment studies two disruptions: stale summaries and per-domain controller failures. Per-domain controller failure is a structurally HiRoute-specific stressor because the centralized-directory baseline does not have per-domain controllers; it appears here only as a non-peer availability reference, not as a peer that competes on this stressor. The goal is therefore not to claim universal robustness superiority but to characterize how HiRoute degrades and recovers under stressors that can arise on a real deployment of the architecture.

\begin{figure}[t]
\centering
\includegraphics[width=0.95\linewidth]{results/figures/mainline/fig_robustness.pdf}
\caption{Degradation profile under controller-down and stale-summary stressors. HiRoute degrades and recovers under per-domain controller failure with a bounded minimum-success floor; the centralized directory is shown as a non-peer availability reference because it does not expose per-domain controllers to this stressor.}
\label{fig:robust}
\end{figure}

Figure~\ref{fig:robust} reports a degradation profile rather than a winner-comparison. On the current workload, HiRoute's controller-down trace shows a measurable success drop with bounded recovery to a minimum-success floor of approximately 0.73, accompanied by additional recovery probes that the manifest-fallback and reliability-cache paths consume to find an alternative answer. Stale summaries do not materially reduce terminal success on either scheme on this workload, which indicates that the staleness stressor is presently weak rather than that the system is uniformly stale-tolerant; both observations should be retested under a clean rerun and a stronger staleness configuration. The 95th-percentile recovery time is computed from one-second time bins, so a value of zero indicates recovery within the first post-event bin rather than the absence of degradation.

\section{Discussion}

\subsection{Positioning as a Routing and Name-Resolution System}

A common criticism of semantic NDN papers is that they sometimes drift away from routing and toward external search. \sysname{} avoids that drift in two ways. First, the final output of discovery is always a standard canonical NDN name, so the data plane remains ordinary NDN retrieval. Second, the inter-domain control plane exports only bounded summaries rather than raw objects. For this reason, the paper is fundamentally about name resolution and routing-state management rather than about replacing NDN with a generic search system.

\subsection{Favorable Operating Regimes}

The design is most likely to help when three conditions hold simultaneously. The first is moderate or strong naming heterogeneity, because exact-name forwarding then becomes brittle. The second is meaningful predicate structure, because the architecture relies on constraints to eliminate large portions of the search space. The third is residual semantic ambiguity within otherwise relevant domains, because the hierarchy and manifest are intended to narrow that ambiguity. If these conditions do not hold, simpler baselines may already be sufficient.

\subsection{Limitations}
\label{sec:limitations}

The design still has limitations. Query decomposition must be available in structured form, whether through an external parser or through offline workload generation. The current prototype instantiates hierarchical summaries as static per-run artifacts rather than as fully online packetized \hslsa{} exchange with sequence numbers and signature validation. The manifest architecture is in place, but the current workload does not yet strongly exercise manifest-rescue behavior because the dominant failure mode remains more often incorrect domain selection than incorrect object selection. Finally, although coarse predicate sketches reveal far less than raw object names, they may still raise privacy concerns in deployments with sensitive geographic or service metadata.

\section{Conclusion}

This paper has presented \sysname{}, a hierarchical constraint-guided name-resolution architecture for NDN-IoT. The design begins from a precise systems claim: IoT queries contain hard constraints and soft semantics, and these two components should not be handled in the same manner. The resulting architecture uses bounded hierarchical summaries, predicate-first elimination, heuristic semantic refinement over admissible cells, manifest-based local resolution, and lightweight reliability feedback. It keeps the forwarding plane unchanged and confines semantic logic to ingress nodes and domain controllers, which makes the design suitable for ndnSIM realization and for incremental deployment arguments. The evaluation is organized around terminal service success, first-fetch correctness, discovery breadth, latency under deadlines, exported-state growth, and degradation under staleness and failures, because these are the dimensions on which a semantic NDN-IoT routing paper must ultimately be judged. Under the current evidence boundary, the strongest support is bounded fallback and state-budget behavior; routing, ablation, deadline, candidate-contraction, and robustness results remain diagnostic until the corresponding clean-promotion and provenance gates pass.

\balance
\begin{thebibliography}{99}

\bibitem{ndn_iot_invited}
W.~Shang, A.~Bannis, T.~Liang, Z.~Wang, Y.~Yu, A.~Afanasyev, J.~Thompson, J.~Burke, B.~Zhang, and L.~Zhang,
``Named data networking of things,''
\emph{Proceedings of the IEEE}, vol.~107, no.~4, pp.~804--820, 2019.

\bibitem{inf_ndn}
G.~M.~Raza, I.~Ullah, M.~S.~U.~Din, M.~A.~U.~Rehman, and B.-S.~Kim,
``INF-NDN IoT: An intelligent naming and forwarding in named data networking for Internet of Things,''
\emph{IEEE Access}, vol.~12, pp.~114319--114337, 2024.

\bibitem{sanr}
G.~M.~Raza, I.~Ullah, M.~N.~Ali, and B.-S.~Kim,
``SANR-CMF: Semantic-aware naming resolution and cache management framework in information-centric network for Internet of Things,''
\emph{IEEE Internet of Things Journal}, vol.~12, no.~17, pp.~35447--35466, 2025.

\bibitem{fuzzy}
K.~Chan, B.~Ko, S.~Mastorakis, A.~Afanasyev, and L.~Zhang,
``Fuzzy Interest forwarding,''
in \emph{Proc. Asian Internet Engineering Conf.}, 2017, pp.~31--37.

\bibitem{semantic_distribution}
G.~M.~Raza, T.~T.~Sari, G.~Secinti, K.~Aurangzeb, M.~N.~Ali, and B.-S.~Kim,
``Semantic-driven intelligent name distribution and retrieval for named data networking,''
in \emph{Proc. IEEE Consumer Communications \& Networking Conf.}, 2025.

\bibitem{context_iot}
M.~Antunes, J.~Quevedo, D.~Gomes, and R.~L.~Aguiar,
``Improve contextual IoT service discovery with semantic models,''
in \emph{Proc. IEEE Global Communications Conference Workshops}, 2020.

\bibitem{routing_scalability}
Y.~Zhang, Z.~Xia, A.~Afanasyev, and L.~Zhang,
``A note on routing scalability in named data networking,''
in \emph{Proc. IEEE Int. Conf. Computer Communications Workshops}, 2018.

\bibitem{scalable_forwarding}
H.~Yuan, T.~Song, and P.~Crowley,
``Scalable NDN forwarding: Concepts, issues and principles,''
in \emph{Proc. IEEE Int. Conf. Network Protocols}, 2012.

\bibitem{odr}
O.~Ascigil, S.~Rene, I.~Psaras, and G.~Pavlou,
``On-demand routing for scalable name-based forwarding,''
in \emph{Proc. ACM Conf. Information-Centric Networking}, 2018, pp.~1--10.

\bibitem{samba}
A.~Esmaeili and A.~Mtibaa,
``SAMBA: Scalable approximate forwarding for NDN implicit FIB aggregation,''
in \emph{Proc. IEEE Int. Conf. Cloud Networking}, 2024.

\bibitem{nlsr}
L.~Wang, V.~Lehman, A.~K.~M.~M.~Hoque, B.~Zhang, Y.~Yu, and L.~Zhang,
``A secure link state routing protocol for NDN,''
\emph{IEEE Access}, vol.~6, pp.~10470--10487, 2018.

\bibitem{istack}
T.~Li, T.~Song, and Y.~Yang,
``iStack: A general and stateful name-based protocol stack for named data networking,''
in \emph{Proc. USENIX NSDI}, 2024, pp.~1037--1053.

\bibitem{hrouting}
V.~Lehman, A.~Gawande, B.~Zhang, L.~Zhang, R.~Aldecoa, D.~Krioukov, and L.~Wang,
``An experimental investigation of hyperbolic routing with a smart forwarding plane in NDN,''
in \emph{Proc. ACM Conf. Information-Centric Networking}, 2016, pp.~71--80.

\end{thebibliography}

\end{document}
